\begin{ouverture}[college1]{3}{États et changements d’état : le modèle}{Partie 1 · Organisation et transformations de la matière}{Que devient l’eau quand elle se vaporise ?}
\objectif{Interpréter les états physiques avec des particules.}
\objectif{Lire une courbe de chauffage et reconnaître un palier.}
\objectif{Expliquer conservation de la masse et variation du volume.}
\prerequis{États solide, liquide, gazeux ; température.}
\end{ouverture}
\pcPage{Le cours}{États et changements d’état : le modèle}
\phantomsection\label{ess-3}
\essSection{1}{Un modèle pour expliquer}
Dans un solide, les particules sont proches et leur position reste organisée ou liée autour de positions moyennes ; elles vibrent. Dans un liquide, elles sont proches mais peuvent se déplacer les unes par rapport aux autres. Dans un gaz, elles sont beaucoup plus espacées et mobiles. Le modèle explique forme et volume sans représenter fidèlement la taille des particules.\par
\pcFigure{figures/c03-etats.pdf}{Les dessins gardent les mêmes particules ; leur disposition change selon l’état.}
\essSection{2}{Nommer les transformations}
Fusion : solide vers liquide ; solidification : liquide vers solide. Vaporisation : liquide vers gaz ; liquéfaction : gaz vers liquide. La vaporisation peut se produire à la surface, par évaporation, sans ébullition. La sublimation et la condensation solide relient directement solide et gaz. Une buée est composée de fines gouttelettes liquides ; la vapeur d’eau gazeuse est invisible.\par
\clearpage
\essSection{3}{Lire une courbe expérimentale}
À pression constante, le changement d’état d’un corps pur se produit à température constante tant que les deux états coexistent : un palier apparaît dans le modèle idéal. Pour l’eau pure sous une pression voisine de $1{,}013\unite{bar}$, fusion à $0\,^{\circ}\mathrm C$ et ébullition vers $100\,^{\circ}\mathrm C$. Une pression différente ou des substances dissoutes modifient ces repères.\par
\pcFigure{figures/c03-courbe.pdf}{Pendant le palier, l’énergie transférée contribue au changement d’état ; la température ne suffit pas à décrire le phénomène.}
\essSection{4}{Ce qui se conserve}
Lors d’un changement d’état sans fuite, la masse se conserve : les molécules ne sont ni détruites ni créées. Le volume peut changer, car leur organisation et leur écartement changent. La glace occupe plus de volume que la même masse d’eau liquide, particularité importante pour les canalisations et les milieux naturels. La dilatation de l’eau liquide et la fonte des glaces continentales contribuent à l’élévation du niveau marin.\par
\begin{pcEncadre}{method}{Un exemple expliqué}Un flacon fermé contient $50\unite g$ d’eau, puis cette eau gèle entièrement. La masse d’eau reste $50\unite g$ ; son volume augmente. Si l’on observe seulement $48\unite g$ après chauffage dans un récipient ouvert, cela ne réfute pas la conservation : une partie de l’eau a pu quitter le récipient à l’état gazeux.\end{pcEncadre}
\pcPage{Expérimenter et retenir}{États et changements d’état : le modèle}
\phantomsection\label{activite-3}
\begin{pcEncadre}{method}{Investigation en classe}\textbf{Question.} Comment évolue la température pendant la fusion ? Le professeur fournit un mélange eau-glace et un dispositif de mesure. Relever la température à intervalles réguliers, sans mettre la sonde contre la paroi. Noter simultanément les états présents. Tracer température en fonction du temps. Distinguer un vrai palier des petites variations de mesure. \textbf{Prolongement documentaire :} comparer glace flottante et glacier continental ; justifier pourquoi leurs fontes n’ont pas le même effet direct sur le niveau de la mer.\end{pcEncadre}
\pcSousTitre{Ce que je dois savoir expliquer}\begin{questions}
\item Un changement d’état conserve les espèces chimiques.
\item Les particules du gaz sont bien plus espacées que celles du liquide.
\item La vapeur d’eau gazeuse est invisible ; la buée est liquide.
\item Les températures de changement d’état dépendent notamment de la pression.
\item Masse conservée ne signifie pas volume conservé.
\end{questions}
\begin{pcEncadre}{attention}{Une erreur à éviter}« L’eau bout toujours à $100\,^{\circ}\mathrm C$. » Cette valeur suppose une eau pure et une pression proche de la pression atmosphérique au niveau de la mer.\end{pcEncadre}
\pcRepere{cor-activite-3}{Pistes pour l’investigation}\pcGrille
\clearpage\markright{Je démarre}
\rubrique{Je démarre}
\par\noindent\begin{minipage}{\linewidth}
\begin{exercice}[1]{Nommer un passage}
\begin{questions}
\item Comment nommer le passage de la glace à l’eau liquide ?
\item Et celui de la vapeur d’eau à l’eau liquide ?
\end{questions}
\end{exercice}
\end{minipage}
\par\noindent\begin{minipage}{\linewidth}
\begin{exercice}[1]{À l’échelle des particules}
\begin{questions}
\item Dans quel état les particules sont-elles les plus éloignées ?
\item Un liquide a-t-il une forme propre ?
\end{questions}
\end{exercice}
\end{minipage}
\par\noindent\begin{minipage}{\linewidth}
\begin{exercice}[1]{Buée et vapeur}
\begin{questions}
\item La buée visible sur une vitre est-elle un gaz ?
\item Pourquoi voit-on un nuage au-dessus d’une casserole ?
\end{questions}
\end{exercice}
\end{minipage}
\par\noindent\begin{minipage}{\linewidth}
\begin{exercice}[1]{Une masse conservée}
\begin{questions}
\item $80\unite g$ de glace fondent dans un récipient fermé. Quelle masse d’eau obtient-on ?
\item Le volume est-il forcément le même ?
\end{questions}
\end{exercice}
\end{minipage}
\par\noindent\begin{minipage}{\linewidth}
\begin{exercice}[1]{Évaporation}
\begin{questions}
\item Un linge sèche sans bouillir. Quel changement d’état a lieu ?
\item Quelle espèce quitte le linge ?
\end{questions}
\end{exercice}
\end{minipage}
\clearpage\markright{Je m’entraîne}
\rubrique{Je m’entraîne}
\par\noindent\begin{minipage}{\linewidth}
\begin{exercice}[2]{Lire un palier}
\pcFigureEx{figures/ex-fusion.pdf}{}
\begin{questions}
\item Une eau pure reste à $0\,^{\circ}\mathrm C$ tandis que la glace fond. Quels états coexistent ?
\item De l’énergie est-elle néanmoins transférée ?
\end{questions}
\end{exercice}
\end{minipage}
\par\noindent\begin{minipage}{\linewidth}
\begin{exercice}[2]{Choisir un état}
\begin{questions}
\item Sous pression usuelle, de l’eau pure à $-10\,^{\circ}\mathrm C$ est-elle normalement solide ou liquide ?
\item Et à $25\,^{\circ}\mathrm C$ ?
\end{questions}
\end{exercice}
\end{minipage}
\par\noindent\begin{minipage}{\linewidth}
\begin{exercice}[2]{Reconnaître les étapes}
\pcFigureEx{figures/ex-ebullition.pdf}{}
\begin{questions}
\item Une courbe monte de $20$ à $100\,^{\circ}\mathrm C$, puis présente un palier. Décrire les deux étapes.
\item Que faut-il connaître pour attribuer le palier à l’eau ?
\end{questions}
\end{exercice}
\end{minipage}
\par\noindent\begin{minipage}{\linewidth}
\begin{exercice}[2]{Un bilan de masse}
\begin{questions}
\item Un récipient ouvert contient $150\unite g$ d’eau avant chauffage, puis $138\unite g$. Quelle masse a quitté le récipient ?
\item S’agit-il nécessairement d’une transformation chimique ?
\end{questions}
\end{exercice}
\end{minipage}
\par\noindent\begin{minipage}{\linewidth}
\begin{exercice}[2]{Comparer deux courbes}
\pcFigureEx{figures/ex-deux-courbes.pdf}{Deux mesures à pression constante ; la précision reste limitée.}
\begin{questions}
\item L’échantillon A présente un palier net ; B change d’état sur une plage de températures. Lequel est le plus compatible avec un corps pur à pression constante ?
\item Cette seule observation prouve-t-elle toute sa composition ?
\end{questions}
\end{exercice}
\end{minipage}
\clearpage\markright{Je raisonne}
\rubrique{Je raisonne}
\par\noindent\begin{minipage}{\linewidth}
\begin{exercice}[3]{Ne pas enfermer un chauffage}
\begin{questions}
\item Pourquoi ne faut-il pas chauffer de l’eau dans un récipient hermétiquement fermé improvisé ?
\item Quel lien avec le volume du gaz ?
\end{questions}
\end{exercice}
\end{minipage}
\par\noindent\begin{minipage}{\linewidth}
\begin{exercice}[3]{La banquise et le glacier}
\begin{questions}
\item Distinguer une glace déjà flottante d’une glace continentale.
\item Pourquoi la fonte d’un glacier continental fait-elle monter le niveau de la mer ?
\end{questions}
\end{exercice}
\end{minipage}
\par\noindent\begin{minipage}{\linewidth}
\begin{exercice}[3]{Critiquer un dessin}
\pcFigureEx{figures/ex-particules-grosses.pdf}{Représentation de l’élève : les particules sont dessinées plus grosses après chauffage.}
\begin{questions}
\item Un élève dessine des molécules plus grosses dans l’eau liquide chaude. Quelle modification proposer ?
\item Que peut changer le chauffage ?
\end{questions}
\end{exercice}
\end{minipage}
\par\noindent\begin{minipage}{\linewidth}
\begin{exercice}[3]{Pression et température}
\begin{questions}
\item Un document annonce une ébullition de l’eau pure à 90 °C en altitude. Est-ce incompatible avec le cours ?
\item Quelle donnée faut-il comparer ?
\end{questions}
\end{exercice}
\end{minipage}
\par\noindent\begin{minipage}{\linewidth}
\begin{exercice}[3]{Construire un argument}
\begin{questions}
\item À $0\,^{\circ}\mathrm C$, deux récipients contiennent l’un de l’eau liquide et l’autre un mélange eau-glace. Peut-on les distinguer par la seule température ?
\item Quelle observation ajouter ?
\end{questions}
\end{exercice}
\end{minipage}
\pcPage{Défi du chapitre}{États et changements d’état : le modèle}
\phantomsection\label{defi-ex-3}
\pcDefiTitre{3}{Le flacon qui semble perdre de la matière}
\begin{questions}
\item Un montage ouvert pèse $312\unite g$ avant et $300\unite g$ après chauffage. Calculer la différence.
\item Proposer une explication compatible avec la conservation de la masse.
\item Comment vérifier l’explication sans fermer dangereusement le montage ?
\end{questions}
\pcGrille
